\ifx\allfiles\undefined
\documentclass[12pt, a4paper, oneside, UTF8]{ctexbook}
\def\path{../config}
\input{\path/package.tex}

% 定理环境设置
\input{\path/theorem1_zh.tex}

\input{\path/custom.tex}

% 修改参数改变封面样式，0 默认原始封面、内置其他1、2、3种封面样式
\def\myIndex{3}


\ifnum\myIndex>0
    \input{\path/cover_package_\myIndex} 
\fi

\def\myTitle{复变函数与积分变换笔记}
\def\myAuthor{M.K.}
\def\myDateCover{\today}
\def\myDateForeword{\today}
\def\myForeword{前言}
\def\myForewordText{
    
    这是一个基于\LaTeX{}模板的复变函数与积分变换笔记．
    内容参考了包括但不限于以下资料：
    \begin{itemize}
        \item 《复变函数与积分变换》（科学出版社）第三版
        \item 《复变函数习题精选精讲》（山东科学技术出版社）第一版
    \end{itemize}

    封面来源于\href{https://www.pixiv.net/users/2642047}{アシマ / Ashima}的作品\href{https://www.pixiv.net/artworks/147092236}{Assembly}
    如有侵权请联系我删除．
}
\def\mySubheading{副标题}


\begin{document}
%\input{../config/cover}
\else
\fi
\chapter{傅里叶变换}

\section{傅里叶积分与傅里叶变换}

\subsection{傅里叶积分定理}

\begin{theorem}[傅里叶积分定理]
设 $f\colon\R\to\C$ 满足：
\begin{enumerate}[label=\roman*.]
    \item Dirichlet 条件：$f$ 在 $\R$ 上分段连续，且在每个有限区间上只有有限个极值点；
    \item 绝对可积性：$\displaystyle\int_{-\infty}^{+\infty}\abs{f(x)}\,\d x < \infty$。
\end{enumerate}
即$f\in L^1(\R)$，则在 $f$ 的连续点 $x$ 处有
\[
\frac{1}{2\pi}\int_{-\infty}^{+\infty}\!\left[\int_{-\infty}^{+\infty} f(t)\e^{-i\omega t}\,\d t\right]\e^{i\omega x}\,\d\omega
= \begin{cases}
    f(x), & f \text{ 在 } x \text{ 处连续} \\[1em]
    \dfrac{f(x^+)+f(x^-)}{2}, & f \text{ 在 } x \text{ 处不连续}
    \end{cases}
\]

\end{theorem}

\subsection{傅里叶变换}

\begin{definition}[积分变换]
    通过积分运算将像原函数 $f(t)$ 映射为像函数 $F(s)$ 的变换称为\textbf{积分变换}，记作
    \[F(s) = \mathcal{T}[f(t)] = \int_{a}^{b} K(s,t)f(t)\,\d t,\]
其中 $K(s,t)$ 称为\textbf{核函数}。
\end{definition}

\begin{definition}[傅里叶变换]
设 $f\in L^1(\R)$，定义 $f$ 的\textbf{傅里叶变换}为
\[
\mathcal{F}[f](\omega) = \hat{f}(\omega) \defeq \int_{-\infty}^{+\infty} f(x)\e^{-i\omega x}\,\d x,\quad \omega\in\R.
\]
\end{definition}

\begin{definition}[傅里叶逆变换]
定义 $f$ 的\textbf{傅里叶逆变换}为
\[
\mathcal{F}^{-1}[F](x) \defeq \frac{1}{2\pi}\int_{-\infty}^{+\infty} F(\omega)\e^{i\omega x}\,\d\omega,\quad x\in\R.
\]
\end{definition}

\section{单位脉冲函数}

\begin{definition}[单位脉冲函数 / Dirac $\delta$ 函数]
    单位脉冲函数 $\delta(t)$ 是满足以下两个条件的\textbf{广义函数}：
    \begin{tasks}[label=\roman*.](2)
        \task $\displaystyle\delta(t) = \begin{cases}0, & t\neq 0\\ \infty, & t = 0\end{cases}$；
        \task $\displaystyle\int_{-\infty}^{+\infty}\delta(t)\,\d t = 1$。
    \end{tasks}
\end{definition}

\begin{theorem}[$\delta$ 函数的筛选性质]
    对任意连续函数 $\varphi$，
    \[
        \int_{-\infty}^{+\infty}\delta(t)\varphi(t)\,\d t = \varphi(0),\qquad
        \int_{-\infty}^{+\infty}\delta(t-t_0)\varphi(t)\,\d t = \varphi(t_0)
    \]
\end{theorem}

\section{傅里叶变换的性质}

\subsection{基本性质}

\begin{theorem}[线性性]
    设 $f,g\in L^1(\R)$，$\alpha,\beta\in\C$。则
    \[
        \mathcal{F}[\alpha f + \beta g] = \alpha\mathcal{F}[f] + \beta\mathcal{F}[g]
    \]
    \[
        \mathcal{F}^{-1}[\alpha F + \beta G] = \alpha\mathcal{F}^{-1}[F] + \beta\mathcal{F}^{-1}[G]
    \]
\end{theorem}

\begin{theorem}[位移性质]
    设 $f\in L^1(\R)$，$t_0\in\R$。则
    \[
        \mathcal{F}[f(t-t_0)](\omega) = \e^{-i\omega t_0}\hat f(\omega),\qquad
        \mathcal{F}[\e^{i\omega_0 t} f(t)](\omega) = \hat f(\omega-\omega_0)
    \]
\end{theorem}

\begin{theorem}[对称性质]
    设 $f\in L^1(\R)$。则
    \[
        \mathcal{F}[f(-t)](\omega) = \hat f(-\omega),\qquad
        \mathcal{F}[\hat f(t)](\omega) = 2\pi f(-\omega)
    \]
\end{theorem}

\begin{theorem}[微分性质]
    设 $f,f'\in L^1(\R)$。则
    \[
        \mathcal{F}[f'(t)] = i\omega \hat f(\omega)
    \]
    若 $f$ 连续，$f^{(k)}\in L^1(\R)$（$1\le k\le n$），$f^{(k)}(t)\to 0$（$\abs{t}\to\infty$），则
    \[
        \mathcal{F}[f^{(n)}(t)] = (i\omega)^n \hat f(\omega)
    \]
\end{theorem}

\begin{theorem}[积分性质]
    设 $g(t) = \displaystyle\int_{-\infty}^{t} f(\tau)\,\d\tau$，且 $\hat f(0)=0$。则
    \[
        \mathcal{F}[g(t)] = \frac{1}{i\omega}\hat f(\omega)
    \]
\end{theorem}

\begin{theorem}[乘 $t^n$ 性质]
    设 $f\in L^1(\R)$，$tf(t)\in L^1(\R)$。则 $\hat f$ 可微且
    \[
        \mathcal{F}[-it f(t)] = \hat f'(\omega)
    \]
\end{theorem}

\subsection{常用函数的傅里叶变换}

\begin{theorem}[常用傅里叶变换对]
设 $a>0$。则
\begin{enumerate}
    \item 单位脉冲：$\mathcal{F}[\delta(t)] = 1$；
    \item 矩形脉冲：$u(t)=\begin{cases}1, & \abs{t}<a\\ 0, & \abs{t}>a\end{cases}$，则$\mathcal{F}[u(t+a)-u(t-a)] = \dfrac{2\sin a\omega}{\omega}$；
    \item 指数衰减：$\mathcal{F}[\e^{-a\abs{t}}] = \dfrac{2a}{a^2+\omega^2}$；
    \item 高斯函数：$\mathcal{F}[\e^{-a t^2}] = \sqrt{\dfrac{\pi}{a}}\e^{-\omega^2/(4a)}$；
    \item 单边指数：$\mathcal{F}[u(t)\e^{-at}] = \dfrac{1}{a+i\omega}$；
    \item 符号函数：$\mathcal{F}[\operatorname{sgn} t] = \dfrac{2}{i\omega}$；
    \item 单位阶跃：$\mathcal{F}[u(t)] = \pi\delta(\omega) + \dfrac{1}{i\omega}$。
\end{enumerate}
\end{theorem}

\section{卷积}

\subsection{卷积的定义}

\begin{definition}[卷积]
设 $f,g\in L^1(\R)$。定义 $f$ 与 $g$ 的\textbf{卷积}为
\[
(f*g)(t) \defeq \int_{-\infty}^{+\infty} f(\tau)g(t-\tau)\,\d\tau,\quad t\in\R.
\]
\end{definition}

\begin{theorem}[卷积的基本性质]
设 $f,g,h\in L^1(\R)$。则
\begin{enumerate}
    \item 交换律：$f*g = g*f$
    \item 结合律：$(f*g)*h = f*(g*h)$
    \item 线性性：$f*(\alpha g + \beta h) = \alpha(f*g) + \beta(f*h)$
    \item 微分：若 $f$ 可微且 $f'\in L^1(\R)$，则 $(f*g)' = f'*g = f*g'$
\end{enumerate}
\end{theorem}

\subsection{卷积定理}

\begin{theorem}[卷积定理（Fourier 变换）]
设 $f,g\in L^1(\R)$。则 $f*g\in L^1(\R)$，且
\[
\mathcal{F}[f*g] = \mathcal{F}[f]\cdot \mathcal{F}[g],\qquad
\mathcal{F}[f\cdot g] = \frac{1}{2\pi}\mathcal{F}[f]*\mathcal{F}[g].
\]
\end{theorem}


\ifx\allfiles\undefined
\end{document}
\fi